\clearpage\fancyhead[L]{\small BEOBACHTUNG.}\hypertarget{native-353-lower}{}\label{native-353-lower}
\begin{center}\textbf{BEOBACHTUNG.}\par\vspace{4pt}\hrule\vspace{4pt}\fontsize{10}{12}\selectfont Göttingische gelehrte Anzeigen. Heft 38, S. 369\;..\;370, 7. März 1812.\par\vspace{4pt}\hrule\vspace{4pt}\end{center}

\par\indent Eine so seltene Beobachtung wie die von \emph{neun} Sternbedeckungen an zwei Abenden verdient es wohl, in diesen Blättern festgehalten zu werden.

\Needspace{27\baselineskip}\par\indent Am 19. Februar dieses Jahres wurden an der hiesigen Sternwarte die folgenden sechs Sternbedeckungen beobachtet.

\begin{center}\fontsize{9}{11}\selectfont\setlength{\tabcolsep}{2.5pt}\renewcommand{\arraystretch}{1.15}\begin{tabular}{lrrrll}
$\gamma$ Tauri, Austritt & $7^{\mathrm u}$ & $22^{\mathrm m}$ & $15^{\mathrm s}0$ & m.\,Z. & Hr. Gerling\\
71 Tauri, Eintritt & $9$ & $56$ & $3.8$ & — & Prof. Harding\\
 & $9$ & $56$ & $3.9$ & — & Prof. Gauss\\
 & $9$ & $56$ & $3.4$ & — & Hr. Gerling\\
$\theta^1$ Tauri, Eintritt & $10$ & $58$ & $56.9$ & — & Prof. Harding\\
 & $10$ & $58$ & $57.4$ & — & Prof. Gauss\\
 & $10$ & $58$ & $57.9$ & — & Hr. Gerling\\
Austritt & $11$ & $55$ & $34.6$ & — & Prof. Gauss\\
 & $11$ & $55$ & $38.6$ & — & Prof. Harding\\
$\theta^2$ Tauri, Eintritt & $10$ & $58$ & $12.0$ & — & Prof. Gauss und Harding zugleich\\
Austritt & $11$ & $59$ & $6.6$ & — & Prof. Gauss\\
 & $11$ & $59$ & $8.6$ & — & Prof. Harding\\
16c Mayer, Eintritt & $12$ & $3$ & $39.5$ & — & Prof. Harding\\
160 Mayer, Eintritt & $12$ & $3$ & $40.0$ & — & Prof. Gauss\\
162 Mayer, Eintritt & $12$ & $26$ & $27.0$ & — & Prof. Gauss\\
 & $12$ & $26$ & $27.5$ & — & Harding\\
\end{tabular}\end{center}

\par\indent Professor \textsc{Gauß} beobachtete mit dem 10-füßigen \textsc{Herschel}-Teleskop, Professor \textsc{Harding} mit einem 5-füßigen Teleskop und Herr Gerling mit dem 5-füßigen Dollond.

\Needspace{13\baselineskip}\par\indent Am 20. Februar wurden noch die folgenden 3 Bedeckungen beobachtet. Dabei benutzte Herr Professor \textsc{Harding} das 5-füßige und Herr \textsc{Gerling} das 10-füßige Teleskop.

\begin{center}\fontsize{9}{11}\selectfont\setlength{\tabcolsep}{2.5pt}\renewcommand{\arraystretch}{1.15}\begin{tabular}{lrrrll}
111 Tauri, Eintritt & $11^{\mathrm u}$ & $7^{\mathrm m}$ & $46^{\mathrm s}8$ & m.\,Z. & Hr. Gerling\\
Anonyma, Eintritt & $11$ & $36$ & $54.0$ & — & Prof. Harding\\
117 Tauri, Eintritt & $13$ & $0$ & $25.0$ & — & Hr. Gerling\\
 & $13$ & $0$ & $25.1$ & — & Prof. Harding\\
\end{tabular}\end{center}
